\documentclass[11pt]{article}
\usepackage[margin=0.8in]{geometry}
\usepackage{enumitem}
\usepackage{setspace}
\setlist[itemize]{leftmargin=1.15em, itemsep=0.1em, topsep=0.18em}
\setlist[enumerate]{leftmargin=1.35em, itemsep=0.1em, topsep=0.18em}
\setlength{\parindent}{0pt}
\setlength{\parskip}{0.26em}
\setstretch{1.05}
\pagestyle{plain}

\begin{document}

\begin{center}
{\Large\bfseries Microeconomic Theory II}\\[0.3em]
Sample syllabus\\[0.15em]
Ke Cheng HSU\\
Department of Economics, National Chengchi University\\
Master's core, 3 credits, second semester, taught in English
\end{center}

\textbf{Place in the curriculum.}
This course continues Microeconomic Theory I and matches the department's Microeconomic Theory (II). Sophomore Microeconomics has already introduced monopoly, oligopoly, and the Edgeworth box. This course requires calculus, mixed-strategy Nash equilibrium, and a statement of the two welfare theorems with the assumption each one uses. An existence proof that uses a fixed-point theorem is left for the advanced graduate course, which uses Mas-Colell, Whinston, and Green.

\textbf{Prerequisite.}
Microeconomic Theory I, or an equivalent course in calculus-based consumer and producer theory. The first class recalls the condition that price equals marginal cost in a competitive industry, because monopoly and oligopoly are departures from it.

\textbf{Textbooks.}
Varian, Hal R., and Marc Melitz. \emph{Intermediate Microeconomics: A Modern Approach}. 10th edition. W.\ W.\ Norton, 2024. The same edition as in Theory I. Games, Chapters 29--30, are taught before oligopoly, Chapter 28, because the oligopoly models use those equilibrium concepts. The book prints Chapter 28 first. Chapter 27, factor markets, and Chapter 31, behavioral economics, are omitted. Auctions, Chapter 18, are taught in Uncertainty and Information Economics. Chapter 37 is limited to the lemons market. Signaling, moral hazard, and screening are taught in that elective.

Steven Tadelis, \emph{Game Theory: An Introduction}, is an optional reference for the game weeks. Students are not asked to buy it. Problem sets follow the exercises in Varian and Melitz, with one extra game written on the sheet when the chapter exercise is thin.

\textbf{How the course runs.}
Six problem sets, due at the start of class. A set from the first half asks for an equilibrium price and a surplus comparison. A set from the second half asks for an Edgeworth-box allocation or a public-good quantity, and for the assumption that makes the welfare theorem go through. Students may discuss the method in groups of up to three and must hand in their own write-up. One class hour every second week is reserved for students to present a solution. Office hours are by appointment and are for a model that has been started.

\textbf{Learning outcomes.}
By the end of the semester, students should be able to
\begin{itemize}
\item derive a monopolist's marginal revenue and solve uniform pricing, third-degree price discrimination, and a two-part tariff,
\item find pure-strategy and mixed-strategy Nash equilibria, and subgame-perfect equilibria in a two-stage game,
\item solve Cournot, Bertrand, and Stackelberg competition and compare the prices,
\item construct an Edgeworth-box equilibrium, state both welfare theorems, and name the assumption that fails under an externality,
\item write the Samuelson condition for a public good.
\end{itemize}

\textbf{Assessment.}
Problem sets 25 percent. Midterm examination 30 percent, covering weeks 1--7. Final examination 45 percent, on this course. A problem may use a method from Theory I inside a question about this semester's material. Exams are closed book.

\textbf{Weekly plan.}
\begin{enumerate}
\item Monopoly pricing and welfare. Chapter 25. Marginal revenue for a linear demand. The quantity where marginal revenue equals marginal cost. The deadweight-loss triangle against the competitive quantity. A markup formula in terms of the elasticity.
\item Price discrimination and a two-part tariff. Chapter 26. Third-degree discrimination across two markets. A two-part tariff that extracts consumer surplus when the monopolist knows the demand. Bundling only as a short example if time remains in the week.
\item Dominant strategies, pure Nash equilibrium, and mixed Nash equilibrium. Chapter 29. Best responses. A coordination game with two pure equilibria. A mixed equilibrium in a $2\times 2$ game, computed by making the other player willing to mix.
\item Extensive-form games and subgame perfection. Selected sections of Chapter 30, Game Applications. Backward induction in a two-stage game. A noncredible threat. One entry game.
\item Cournot and Bertrand competition. Chapter 28. Reaction functions in quantities. The Bertrand paradox with homogeneous goods and constant marginal cost. Differentiated Bertrand as the case in which price stays above cost.
\item Stackelberg competition. A second week on Chapter 28. The leader's quantity, the follower's reaction, and why the leader's profit exceeds the Cournot profit.
\item Collusion with a trigger strategy. Lecture notes. An infinitely repeated Bertrand game. The lowest discount factor that sustains the monopoly price. This model is shorter than a chapter of the textbook.
\item Midterm examination on weeks 1--7.
\item Pure exchange and the Edgeworth box. Chapter 32. The contract curve. A competitive equilibrium as a price line through the endowment that supports a Pareto allocation. Excess demand.
\item General equilibrium with production. Chapter 33. A Robinson Crusoe economy. The production-possibility frontier. The condition that the marginal rate of substitution equals the marginal rate of transformation.
\item The two welfare theorems. Chapter 34. The first theorem under local nonsatiation and price-taking. The second theorem under convex preferences and convex technology, with a transfer of endowments. No fixed-point argument for existence.
\item Externalities. Chapter 35. A consumption externality in an Edgeworth box, where the competitive allocation misses the contract curve. A Pigouvian tax. The Coase theorem when property rights are clear and bargaining has no cost.
\item Public goods and the Samuelson condition. Chapter 36. The sum of marginal rates of substitution equals marginal cost. Free riding in a voluntary-contribution game. One mechanism that asks each person to report a value, stated without a general proof.
\item The lemons market. Selected sections of Chapter 37. Quality known only to the seller. The average quality that buyers are willing to pay for, and the unraveling when that average falls. Signaling and moral hazard are taught in Uncertainty and Information Economics.
\item Information technology. Chapter 38. Network externalities, switching costs, and systems competition, from the textbook models in that chapter.
\item Workshop. One oligopoly model from weeks 5--6 and one Edgeworth box from week 9, worked in class from a blank page.
\item Review of the welfare theorems and of the market failures in weeks 12--14.
\item Final examination on this course.
\end{enumerate}

\end{document}
